\documentclass[11pt]{article}
\usepackage[margin=1in]{geometry}
\usepackage[utf8]{inputenc}
\usepackage{amsmath,amssymb}
\usepackage{booktabs}
\usepackage{graphicx}
\usepackage{hyperref}
\usepackage{setspace}
\usepackage{natbib}
\hypersetup{colorlinks=true, linkcolor=black, urlcolor=blue, citecolor=black}
\setstretch{1.15}

\title{Transmission Bottlenecks and Battery Storage Mitigation for Spatially Clustered EV Charging: A California Case Study}
\author{Jiewei Grant Li$^{1,*}$, Alan Jenn$^{1}$\\
\small $^1$Institute of Transportation Studies, University of California, Davis\\
\small $^*$Corresponding author: jwgli@ucdavis.edu}
\date{}

\begin{document}
\maketitle

\begin{abstract}
Electric vehicle (EV) charging can change consequential power-sector emissions when spatial clustering interacts with constrained delivery of remote clean generation. We develop a two-layer EV--grid framework for California: (i)~a multi-day, SOC-aware charging demand model that allocates fleet load to transportation analysis zones (TAZs), and (ii)~a Western Electricity Coordinating Council (WECC) Grid Optimized Operation Dispatch (GOOD) model nested with a California meso-grid of aggregated substation hubs. Using PG\&E Grid Resource Integration Portal (GRIP) substations and transmission lines, TAZ centroids are mapped to nearest substations, clustered into 80 meso hubs, and linked to parent balancing-area nodes via capacity-constrained interfaces. Four seasonal weeks and four counterfactuals (S0--S3) isolate a congestion-emissions penalty \(P_{cong}\) and a BESS mitigation metric \(M_{BESS}\). Seasonal-week solves on a local workstation produce structured diurnal EV demand (peak \(\sim\)3.5--3.8~GW for a 2.5~million vehicle fleet) and positive \(P_{cong}\) in all four seasons, motivating storage and transfer upgrades as distinct levers from temporal load shifting alone.
\end{abstract}

\section{Introduction}

Electric vehicles can reduce direct transportation emissions, but charging demand must be delivered through a power system with finite generation, storage, and transmission capacity. Well-to-wheel and life-cycle studies show that EV climate benefits depend on upstream electricity supply rather than tailpipe emissions alone \citep{hawkins_comparative_2013,tamayao_regional_2015,challa_well--wheel_2022}. Average grid-emission factors are insufficient when charging load changes marginal generation, renewable curtailment, and power transfers across constrained networks \citep{ambrose_trends_2020,burton_data-driven_2023}.

A balancing-area consequential-emissions framework can identify how EV charging changes generation, storage, and inter-regional exchange. Aggregate regional modeling can hide an important operational mechanism: EV charging is spatially clustered, while clean generation is often remote. When transmission corridors or local delivery interfaces bind, low-carbon imports can be curtailed and local fossil generation may be dispatched out of merit \citep{bird_wind_2016,ahmed_grid_2020}. This study quantifies how much of the consequential emissions from EV charging is attributable to transmission bottlenecks, and whether strategically located battery energy storage systems (BESS) can reduce congestion-related emissions.

\section{Background and contribution}

Prior EV--grid studies increasingly combine empirical charging profiles with dispatch or capacity-expansion models to estimate marginal power-sector impacts \citep{powell_scalable_2022,jenn_environmental_2020,jenn_emissions_2023,bruchon_cleaning_2024}. Storage has also been proposed as a mitigation resource for charging-related grid stress because it can shift renewable generation into high-demand charging periods and reduce dependence on constrained imports \citep{mowry_grid_2021,nagarajan_value_2018,ziegler_storage_2019,sihvonen_techno-socio-economic_2025}. Yet the emissions value of storage depends on where it is connected and whether it relieves a binding delivery constraint rather than merely shifting energy temporally.

Our contribution is a two-layer EV--grid modeling framework. The first layer uses balancing-area GOOD to estimate consequential generation and capacity expansion across WECC. The second layer allocates California charging demand to TAZs, maps TAZ centroids to nearby substations (PG\&E GRIP plus gateway proxies outside the PGE footprint), and represents a capacity-constrained California meso delivery layer nested inside WECC. We evaluate S0 (no EV), S1 (EV, constrained delivery, no new BESS), S2 (EV, constrained, endogenous BESS), and S3 (EV, relaxed delivery). The key outputs are
\[
P_{cong}=(E_{S1}-E_{S0})-(E_{S3}-E_{S0}),
\qquad
M_{BESS}=(E_{S1}-E_{S0})-(E_{S2}-E_{S0}).
\]

\section{Data and methods}

\subsection{EV charging demand}

The EV module builds from a multi-day, SOC-aware charging simulation with heterogeneous battery capacities, daily VMT, home/work charging access, and charging-frequency preferences. Energy demand accumulates over multiple days; charge events are triggered when accumulated demand approaches usable capacity or a preferred interval is reached. Each event is matched to an empirical charging-session pool by location type, charger level or housing type, and energy bin.

The California aggregate case uses approximately 2.5~million light-duty EVs, 35~miles/day, 0.3~kWh/mile, and charging-energy shares of 68\% residential L1/L2, 4\% workplace L2, 8\% public L2, and 20\% DC fast charging \citep{nrel_2030charging}. Annual TAZ energy totals are combined with the fleet diurnal shape to form hourly TAZ loads without storing full session tables:
\[
L^{EV}_{i,t}=L^{\mathrm{fleet}}_{t}\cdot\frac{w_i}{\sum_j w_j},
\]
where \(w_i\) is annual TAZ charging energy. Home charging is assigned to the vehicle's home TAZ; work and public charging use destination proxies until trip-level destinations are fully integrated.

\begin{figure}[htbp]
  \centering
  \includegraphics[width=0.85\linewidth]{image/ev_load_daily_curve_by_type.png}
  \caption{Empirical-session-based EV charging profiles by charger category and day type (extended abstract). These profiles provide the temporal load shape for the transmission-constrained GOOD model.}
  \label{fig:ev-load}
\end{figure}

\subsection{TAZ--substation mapping and meso aggregation}

Each TAZ \(i\) is represented by its centroid \(c_i\) and assigned to the nearest candidate substation \(s\):
\[
s(i)=\arg\min_{s\in S} dist(c_i,c_s).
\]
Candidate set \(S\) combines PG\&E GRIP \texttt{EDSubstations} (\(n=704\)) with balancing-area gateway points for non-PGE California when statewide CEC/HIFLD downloads are unavailable. Hourly substation load is
\[
L^{EV}_{s,t}=\sum_{i:s(i)=s} L^{EV}_{i,t}.
\]
Substations are clustered into \(N=80\) meso hubs. GRIP \texttt{TransmissionLines} provide hub--hub corridors with transfer capacities inferred from rated kV. Each hub links to a parent California balancing-area node through a bidirectional interface.

\begin{figure}[htbp]
  \centering
  \includegraphics[width=0.72\linewidth]{../../figures/ASTR_diagnostics/substation_ev_concentration.png}
  \caption{EV charging concentration at mapped substations (mean seasonal-week energy). PG\&E GRIP stations shown; southern load is partly represented via BA gateway proxies pending CEC layer ingest.}
  \label{fig:sub-conc}
\end{figure}

\subsection{GOOD and nested delivery layer}

GOOD minimizes operational costs, capital investment, and soft-constraint penalties subject to regional power balance, policy targets, producer limits, load profiles, transmission limits, and storage inventory. We retain the existing WECC balancing-area graph for generation and base load, and add meso hubs that carry only EV Load assets (and optional endogenous Store assets in S2). Nodal balance at meso hubs therefore requires transfers across interfaces and corridors---the congestion mechanism of interest.

Implemented network physics use capacity-constrained directed transfers (GOOD \texttt{Transmission}), a computationally staged surrogate for the DC-flow equations
\[
f_{\ell,t}=b_{\ell}(\theta_{i,t}-\theta_{j,t}),
\qquad
-\bar F_{\ell}-z^T_{\ell}\le f_{\ell,t}\le \bar F_{\ell}+z^T_{\ell}
\]
stated in the extended abstract. Full angle-based DC OPF remains future work.

\subsection{Scenarios and experiment design}

\begin{center}
\begin{tabular}{ll}
\toprule
Scenario & Description \\
\midrule
S0 & Baseline WECC week, no meso EV \\
S1 & EV on meso hubs; tight meso/interface caps; no new BESS \\
S2 & S1 + endogenous BESS on hubs \\
S3 & EV on meso hubs; relaxed (10\(\times\)) meso/interface caps \\
\bottomrule
\end{tabular}
\end{center}

Representative weeks are March, June, September, and December (168~h each), matching prior balancing-area GOOD seasonal experiments. Reporting metrics include top hubs by EV energy, corridor capacities, binding transfers (where available), fossil generation deltas, BESS build, emissions per kWh EV, \(P_{cong}\), and \(M_{BESS}\).

\section{Results}

\subsection{Spatial EV demand}

Fleet seasonal-week peaks are approximately 3.4--3.8~GW (Figure~\ref{fig:ev-weeks}). Aggregation to 80 hubs preserves the system peak while concentrating load on a small set of hubs (Figure~\ref{fig:top-hubs}).

\begin{figure}[htbp]
  \centering
  \includegraphics[width=0.95\linewidth]{../../figures/ASTR_diagnostics/ev_load_seasonal_weeks.png}
  \caption{Aggregate California EV charging load by seasonal week (GW).}
  \label{fig:ev-weeks}
\end{figure}

\begin{figure}[htbp]
  \centering
  \includegraphics[width=0.85\linewidth]{../../figures/ASTR_diagnostics/top20_meso_hubs.png}
  \caption{Top 20 meso hubs by mean seasonal-week EV energy.}
  \label{fig:top-hubs}
\end{figure}

\subsection{Seasonal nested solves}

On a local workstation (\(\sim\)32~GB RAM, Gurobi~13), nested S0--S3 weeks solve for all four seasons. Tables~\ref{tab:june}--\ref{tab:sep-dec} report objective values and CO\(_2\) totals from plant-level generation. Across seasons, relaxed delivery (S3) yields substantially lower CO\(_2\) than constrained EV delivery (S1), so \(P_{cong}>0\). Endogenous BESS on hubs (S2) does not yet show a stable positive \(M_{BESS}\) in these first passes, indicating that storage value is sensitive to shortfall penalties, interface sizing, and BESS cost/siting---a diagnostic result rather than a final policy estimate. For tractability, S2 sites BESS on the top five EV hubs.

\begin{table}[htbp]
\centering
\caption{June representative week: nested GOOD outcomes.}
\label{tab:june}
\begin{tabular}{lrrr}
\toprule
Scenario & Objective & CO\(_2\) (kg) & Nodes \\
\midrule
S0 & \(1.46\times10^8\) & \(1.64\times10^9\) & 16 \\
S1 & \(1.01\times10^8\) & \(1.46\times10^9\) & 96 \\
S2 & \(1.39\times10^8\) & \(1.67\times10^9\) & 96 \\
S3 & \(6.24\times10^7\) & \(6.64\times10^8\) & 96 \\
\midrule
\(P_{cong}\) & & \(8.00\times10^8\) & \\
\(M_{BESS}\) & & \(-2.07\times10^8\) & \\
\bottomrule
\end{tabular}
\end{table}

\begin{table}[htbp]
\centering
\caption{March representative week: nested GOOD outcomes (higher shortfall penalty).}
\label{tab:march}
\begin{tabular}{lrrr}
\toprule
Scenario & Objective & CO\(_2\) (kg) & Nodes \\
\midrule
S0 & \(1.28\times10^8\) & \(2.19\times10^9\) & 16 \\
S1 & \(1.01\times10^8\) & \(1.87\times10^9\) & 96 \\
S2 & \(1.35\times10^8\) & \(2.21\times10^9\) & 96 \\
S3 & \(6.25\times10^7\) & \(6.90\times10^8\) & 96 \\
\midrule
\(P_{cong}\) & & \(1.18\times10^9\) & \\
\(M_{BESS}\) & & \(-3.42\times10^8\) & \\
\bottomrule
\end{tabular}
\end{table}

\begin{table}[htbp]
\centering
\caption{September and December weeks: nested GOOD outcomes.}
\label{tab:sep-dec}
\begin{tabular}{llrrr}
\toprule
Season & Scenario & Objective & CO\(_2\) (kg) & Nodes \\
\midrule
Sep & S0 & \(1.76\times10^8\) & \(4.42\times10^9\) & 16 \\
Sep & S1 & \(1.50\times10^8\) & \(4.14\times10^9\) & 96 \\
Sep & S2 & \(1.64\times10^8\) & \(4.30\times10^9\) & 96 \\
Sep & S3 & \(8.42\times10^7\) & \(1.90\times10^9\) & 96 \\
Dec & S0 & \(1.46\times10^8\) & \(2.17\times10^9\) & 16 \\
Dec & S1 & \(1.22\times10^8\) & \(1.80\times10^9\) & 96 \\
Dec & S2 & \(1.35\times10^8\) & \(1.91\times10^9\) & 96 \\
Dec & S3 & \(7.53\times10^7\) & \(7.49\times10^8\) & 96 \\
\bottomrule
\end{tabular}
\end{table}

\begin{figure}[htbp]
  \centering
  \includegraphics[width=0.75\linewidth]{../../figures/ASTR_diagnostics/P_cong_M_BESS_by_season.png}
  \caption{Congestion-emissions penalty and BESS mitigation by season.}
  \label{fig:pcong}
\end{figure}

\subsection{Computational notes}

Balancing-area weeks remain tractable locally (\(\sim\)7--9~minutes per solve). Nested models grow to \(\sim\)1.85~M rows / 0.73~M columns and typically solve in \(\sim\)8--15~minutes with Gurobi dual simplex and MPS I/O; S2 with storage is restricted to the top five EV hubs with a longer time limit. Full-year nested chronology is not attempted on 32~GB RAM; seasonal weeks are the designed screening design. This machine (i5-14600KF, \(\sim\)32~GB RAM, Gurobi~13) can run the staged seasonal experiment locally.

\section{Discussion}

The framework distinguishes temporal emissions from delivery-constrained emissions. If BESS reduces emissions primarily by charging in renewable-rich hours and discharging during evening EV demand, storage value is temporal. If BESS reduces the gap between constrained and relaxed-transmission cases, it directly mitigates congestion-related emissions. All four seasons show a large S1--S3 gap (\(P_{cong}>0\)). \(M_{BESS}\) is not yet stably positive, consistent with the need to (i)~calibrate unmet-load penalties, (ii)~refine interface capacities from substation load profiles / ICA, and (iii)~site BESS only at high-EV hubs.

Limitations include: nearest-substation screening (not feeder polygons); BA gateway proxies where CEC/HIFLD downloads fail; capacity-constrained transfers rather than DC angles; and proxy work/public destination allocation. SCE/SDGE feeder GIS remain required for the separate distribution-overload pipeline.

\section{Conclusion}

We operationalize the ASTR2026 extended abstract as a runnable TAZ\(\to\)substation\(\to\)meso\(\to\)GOOD pipeline with seasonal S0--S3 counterfactuals. Nested seasonal-week solves demonstrate computational feasibility on a workstation and a positive congestion-emissions penalty under relaxed versus constrained California delivery. Parameter calibration of shortfall costs, interface ratings, and BESS siting will refine \(P_{cong}\) and \(M_{BESS}\) for policy interpretation.

\section*{Data and code}

Scripts \texttt{08\_01}--\texttt{10\_04} live under \texttt{Distribution-Grid-EV-CA/}. GRIP shapefiles are used from
\texttt{data/GRIP\_SHP/.../GRIP\_SHP/}. Outputs: \texttt{data/meso/}, \texttt{astr\_meso\_results/}, \texttt{figures/ASTR\_diagnostics/}.

\newpage
\bibliographystyle{ieeetr}
\bibliography{ch3_transmission_refs}

\end{document}
